\infographicchapter{Indocyanine Green Angiography}{icga.png}
\label{ch:infographic-icga}

\section{Overview}

Indocyanine green angiography (ICGA) is a specialized eye test that uses a safe green dye and gentle near-infrared light to photograph the deep blood vessels at the back of the eye. This deep layer of vessels, called the choroid, sits right behind the retina and supplies it with oxygen and nutrients. Because near-infrared light passes easily through natural eye pigment and thin layers of blood or fluid, ICGA reveals deep problems that standard eye cameras cannot clearly see.

Once injected into an arm vein, almost all of the dye (about 98\%) binds tightly to proteins in your blood. This keeps the dye inside the deep blood vessels longer instead of leaking out quickly, giving eye doctors a crisp map of deep circulation. Doctors most often order ICGA to diagnose conditions like polypoidal choroidal vasculopathy (a type of abnormal bulging blood vessel under the retina), hidden blood vessel growth in wet age-related macular degeneration, chronic fluid buildup under the retina (central serous chorioretinopathy), deep eye inflammation, and certain eye tumors.

ICGA works hand in hand with a more common dye test called fluorescein angiography, which shows the surface vessels of the retina. While ICGA is painless and uses zero X-ray radiation, it does require a quick IV injection and takes longer than a standard eye photo because late pictures are captured up to 30 or 40 minutes afterward. Because the dye is cleared by the liver and standard formulas contain a small amount of sodium iodide, your care team will always check for iodide or dye allergies, thyroid problems, severe liver disease, or pregnancy before the test.

\section{Underlying Physical Principles}

\begin{itemize}
    \item \textbf{Gentle Near-Infrared Light:} A specialized eye camera shines dim, safe near-infrared light into the dilated eye. This light excites the indocyanine green dye traveling through the deep blood vessels, causing the dye to glow softly in the infrared range.
    \item \textbf{Seeing Beneath the Surface:} Unlike regular visible light, near-infrared light slips past natural eye pigments (like melanin) and thin layers of blood or fluid. Because about 98\% of the dye stays attached to blood proteins, it does not leak away quickly, allowing the camera to outline deep choroidal vessels sharply.
    \item \textbf{Sensitive Infrared Detectors:} Indocyanine green glows much more faintly than standard fluorescein dye. To capture clear pictures, the camera uses sensitive infrared detectors and special optical filters that block glare and out-of-focus light.
    \item \textbf{Timed Multi-Phase Pictures:} Taking photos at set times---early (0 to 2 minutes), middle (5 to 15 minutes), and late (20 to 40 minutes)---shows how blood enters, fills, and drains from the back of the eye, highlighting abnormal vessels, leaks, or inflamed spots.
\end{itemize}

\section{Imaging Mechanism}

\begin{enumerate}
    \item \textbf{Eye Alignment and Focus:} You sit comfortably at a special infrared eye camera with your pupils widened by eye drops. The photographer lines up your eye and focuses on the deep layers behind the retina, sometimes setting up the camera to take standard fluorescein photos at the same time.
    \item \textbf{Dye Injection and Timed Photos:} A small dose of indocyanine green dye is injected quickly into an arm vein. As the dye reaches your eye seconds later, the camera captures rapid early images for the first 1 to 2 minutes, followed by scheduled middle- and late-phase photos over the next 30 to 40 minutes.
    \item \textbf{Capturing the Faint Glow:} Special filters block stray reflected light while sensitive infrared sensors pick up the faint glow from the dye and turn it into detailed black-and-white digital images.
    \item \textbf{Eye-Tracking and Image Sharpening:} Built-in eye-tracking software follows tiny natural eye movements and blends several frames together, reducing blur and making deep blood vessels stand out clearly.
    \item \textbf{Side-by-Side Review:} The time-stamped early, middle, and late images appear on the eye specialist's screen, often matched directly with other eye scans like optical coherence tomography (OCT), to guide precise treatment.
\end{enumerate}

\section{Key Technical Parameters}

\noindent\begin{tabularx}{\textwidth}{@{}l X@{}}
\toprule
\textbf{Parameter} & \textbf{Typical Value / Description} \\
\midrule
Physical Principle & Near-infrared light makes protein-bound indocyanine green dye glow inside the deep blood vessels behind the retina \\
Energy / Signal Source & Near-infrared laser eye scanner or filtered infrared eye camera \\
Typical Frequency / Energy & Near-infrared light (around 790--805 nanometers entering, 835 nanometers glowing back) \\
Spatial Resolution & Very fine detail (10--20 micrometers) across standard or wide-angle views of the back of the eye \\
Typical Scan Duration & 30--45 minutes of spaced photos (plus 20--30 minutes beforehand for eye-drop dilation) \\
Contrast Agent & Intravenous indocyanine green dye (usually 25 mg; standard formulas contain up to 5\% sodium iodide) \\
Ionizing Radiation & None (uses safe near-infrared light with zero X-ray radiation) \\
\bottomrule
\end{tabularx}

\section{Clinical Usages}

\subsection{Primary Usages}

\begin{itemize}
    \item \textbf{Polypoidal Choroidal Vasculopathy (PCV):} The most reliable test for spotting tiny, bulging abnormal blood vessels (polyps) and branching vessel networks tucked under the retina, which can cause sudden bleeding or fluid buildup.
    \item \textbf{Wet Age-Related Macular Degeneration (AMD):} Helps locate hidden abnormal blood vessels growing beneath the retina, especially when bleeding or fluid blocks standard visible-light cameras.
    \item \textbf{Central Serous Chorioretinopathy (CSCR):} Pinpoints leaky, swollen deep blood vessels in chronic fluid buildup under the retina, helping doctors target gentle laser treatments (photodynamic therapy) accurately.
    \item \textbf{Deep Eye Inflammation:} Reveals hidden pockets of inflammation or poor blood flow in the deep choroid layer, such as in white dot syndromes and Vogt-Koyanagi-Harada disease.
\end{itemize}

\subsection{Secondary Usages}

\begin{itemize}
    \item \textbf{Evaluating Eye Tumors:} Shows the blood supply inside deep eye growths---such as benign blood vessel tumors (hemangiomas)---and helps tell them apart from cancerous growths.
    \item \textbf{Checking Poor Treatment Response:} Helps re-evaluate patients with wet macular degeneration who are not improving as expected with eye injections, often uncovering hidden PCV or chronic CSCR.
    \item \textbf{Severe Nearsightedness and Tissue Cracks:} Detects subtle deep cracks in the back of the eye (lacquer cracks or angioid streaks) and new abnormal vessels hidden beneath pigment or blood.
\end{itemize}

\section{Advantages}

\begin{itemize}
    \item \textbf{Sees Deep Behind the Retina:} Near-infrared light travels cleanly through natural eye pigment and thin layers of blood or fluid, showing deep problems that standard eye tests miss.
    \item \textbf{Sharp Views of Deep Vessels:} Because about 98\% of the dye binds to blood proteins, it stays inside the vessels longer and outlines deep blood vessel networks clearly.
    \item \textbf{Best Test for Hidden Vessel Bulges:} Uniquely spots polypoidal choroidal vasculopathy and hidden abnormal blood vessels that often stay invisible on regular fluorescein angiography.
    \item \textbf{Reveals Hidden Inflammation and Growths:} Detects deep inflammatory spots and distinctive blood-flow patterns in eye tumors that cannot be seen during a routine eye exam.
    \item \textbf{Zero X-Ray Radiation:} Uses only gentle near-infrared light and can often be done in the same sitting as fluorescein angiography on dual-camera systems.
\end{itemize}

\section{Disadvantages}

\begin{itemize}
    \item \textbf{Requires an IV Injection and Carries Allergy Risk:} Because dye is injected into an arm vein, there is a rare risk of allergic reactions (such as hives, low blood pressure, wheezing, or severe anaphylaxis) as well as mild soreness if dye leaks around the vein.
    \item \textbf{Iodide and Liver Precautions:} Standard ICG dye contains up to 5\% sodium iodide and is cleared by the liver. It is avoided in people with known iodide or ICG allergies and used with caution---or skipped---in people with an overactive thyroid, severe liver disease, severe kidney failure, or pregnancy.
    \item \textbf{Longer Visit and Specialized Equipment:} Because the infrared glow is faint, clinics need specialized infrared cameras and must take spaced late pictures for up to 30 to 40 minutes.
    \item \textbf{Less Detail on Surface Retinal Vessels:} ICGA does not show tiny surface capillaries of the retina as sharply as fluorescein angiography, so doctors often pair the two tests together.
    \item \textbf{Blocked by Dense Cloudiness:} Thick cataracts, heavy bleeding inside the eye, or very thick blood under the retina can still block near-infrared light, and flat 2D photos do not show tissue depth without an OCT scan.
\end{itemize}

\section{Contraindications}

\subsection{Absolute}

\begin{itemize}
    \item \textbf{Severe Allergy to ICG or Iodide:} Anyone who has had a severe allergic reaction (anaphylaxis) to indocyanine green dye or sodium iodide should not receive standard ICG dye. (Iodide-free versions exist only in very limited regions.)
\end{itemize}

\subsection{Relative}

\begin{itemize}
    \item \textbf{Thyroid, Liver, and Kidney Conditions:} Uncontrolled overactive thyroid (hyperthyroidism) or thyroid nodules (since iodide can trigger thyroid flare-ups), severe liver disease (which slows dye clearance), and advanced kidney failure (uremia) require careful medical review.
    \item \textbf{Pregnancy and Eye Conditions:} Elective ICGA is usually postponed during pregnancy unless truly necessary. Caution is also needed for eyes with narrow drainage angles (risk of sudden eye pressure rise from dilating drops) or dense cataracts and heavy eye bleeding that block the camera view.
\end{itemize}

\section{Patient Safety Guidelines}

\begin{itemize}
    \item \textbf{Screening for Allergies, Thyroid, and Liver Issues:} Your care team will ask carefully about any past dye or iodide allergies, overactive thyroid, thyroid nodules, severe liver disease, or severe kidney problems before giving the dye.
    \item \textbf{Pregnancy Precautions:} Patients who could be pregnant are screened beforehand, and routine ICGA is generally delayed during pregnancy unless the medical benefit clearly outweighs any potential risk.
    \item \textbf{Safe IV Placement:} A nurse or technician places a secure arm IV and flushes it gently with saline first to make sure the vein is open, watching the site closely during injection to prevent painful dye leakage under the skin.
    \item \textbf{Emergency Readiness:} Clinics keep emergency allergy medications (such as epinephrine, antihistamines, steroids, oxygen, and IV fluids) right at hand with trained staff, and watch you during and after the test for hives, dizziness, or breathing trouble.
    \item \textbf{Eye-Drop and Thyroid-Test Safety:} Your doctor checks your eye's drainage angle before using dilating drops. Also remember that the small amount of iodide in the dye can interfere with radioactive iodine thyroid tests for at least one week afterward.
\end{itemize}

\section{Prerequisite Tests}

\begin{itemize}
    \item \textbf{Targeted Health Checks (No Routine Blood Work):} Healthy patients usually do not need blood tests beforehand. However, if you have a possible thyroid condition or severe liver disease, your doctor may check thyroid (TSH, free T4) or liver blood tests first.
    \item \textbf{Allergy and Pregnancy Review:} Your team will review your history of iodide or dye allergies, severe asthma, advanced kidney disease, and possible pregnancy. Routine allergy skin-prick testing is not needed.
    \item \textbf{Baseline Eye Exam:} Before dilating your pupils, the clinic checks your vision, eye pressure, and the front chamber of your eye to make sure dilating drops are safe and that your eye is clear enough for good photos.
    \item \textbf{Non-Invasive Retinal Scans First:} Standard color eye photos and a painless optical coherence tomography (OCT) scan---often along with fluorescein angiography---are usually done first to confirm that deep choroidal imaging is needed.
\end{itemize}

\section{Preparing for the Scan}

\begin{itemize}
    \item \textbf{Dilating Drops and Sunglasses:} Expect dilating eye drops 20 to 30 minutes before the photos begin so the camera can see clearly into the back of your eye. Bring dark sunglasses to keep your eyes comfortable in bright light afterward.
    \item \textbf{Bring a Driver and Plan Extra Time:} Arrange for someone to drive you home, as your near vision will be blurry and sensitive to glare for 4 to 6 hours. Plan for a 45- to 60-minute appointment because late photos are taken up to 30 to 40 minutes after the injection.
    \item \textbf{Eating, Medications, and Clothing:} Eat a light meal beforehand---no fasting is required. Take your regular medicines and eye drops as usual, remove contact lenses before dilation, and wear a loose short-sleeved shirt for easy IV access.
    \item \textbf{Share Key Health Details:} Before signing consent, tell your team about any iodide or dye allergy, overactive thyroid, liver disease, pregnancy, or upcoming radioactive iodine thyroid scan.
\end{itemize}

\section{During Scan}

\begin{itemize}
    \item \textbf{Sitting at the Camera:} You will sit upright with your chin resting on a padded chin rest and your forehead gently against a support bar. Simply keep your head steady and look at the small target light inside the lens as instructed.
    \item \textbf{Dye Injection and Spaced Photos:} After a few baseline pictures, the dye is injected into your arm vein. The photographer takes quick photos during the first 1 to 2 minutes, then invites you to sit back and rest between shorter photo sessions at about 5, 15, and 30 minutes.
    \item \textbf{What It Feels and Looks Like:} Because ICGA uses dim near-infrared light, you will see only a soft, faint red glow rather than harsh bright flashes (unless standard fluorescein photos are taken at the same time). You may feel a brief cool feeling in your arm or mild, passing nausea, though nausea is usually much milder than with fluorescein dye.
    \item \textbf{When to Speak Up:} Tell the staff right away if you notice itching, hives, flushing, dizziness, throat tightness, shortness of breath, or burning and swelling at the IV site.
    \item \textbf{How Long It Takes:} Once your pupils are dilated, the photo session itself lasts about 30 to 45 minutes total, with comfortable breaks in between the timed pictures.
\end{itemize}

\section{Normal Appearance}

\begin{itemize}
    \item \textbf{Early Phase (0 to 2 Minutes):} Large deep blood vessels behind the retina fill smoothly and rapidly with bright dye, while surface retinal vessels show up cleanly with no leaking.
    \item \textbf{Middle Phase (5 to 15 Minutes):} The deep tissue layer glows with a soft, even background light as the dye spreads steadily and the optic nerve disc looks gently dark.
    \item \textbf{Late Phase (20 to 40 Minutes):} The background glow fades smoothly and evenly, large deep vessels appear as clean dark outlines against the faint background, the optic disc stays dark, and there are no bright abnormal spots, bulges, or leaky patches.
\end{itemize}

\section{Observing for Abnormal Patterns}

\begin{itemize}
    \item \textbf{Bulging Vessel Spots (PCV):} Small, round, intensely bright spots (polyps) showing up in the first 5 to 6 minutes---often clustered at the edge of a branching web of abnormal vessels---are the hallmark sign of polypoidal choroidal vasculopathy.
    \item \textbf{Hidden Neovascular Plaques and Hot Spots:} Clearly outlined bright spots or broader glowing patches that linger 15 to 30 minutes after injection reveal hidden abnormal blood vessels growing beneath the retina in wet macular degeneration.
    \item \textbf{Leaky, Swollen Deep Vessels:} Widened, winding deep veins surrounded by patchy glowing areas in the middle and late phases point to chronic central serous chorioretinopathy and help guide gentle targeted laser therapy.
    \item \textbf{Dark Spots of Inflammation or Poor Flow:} Evenly scattered round dark dots that persist into the late phase suggest deep inflammatory nodules (granulomas), while larger map-like dark patches indicate areas where deep blood flow is blocked in inflammatory eye conditions.
    \item \textbf{Tumor Patterns and Look-Alikes:} Benign blood vessel tumors (hemangiomas) light up rapidly early on and then quickly fade dark (``wash out'') later. Thick blood, dense pigment, or cataracts can block light and create dark shadows, while harmless fluid pockets can sometimes glow late and mimic abnormal vessels.
\end{itemize}

\section{Dealing with Uncertainties}

\begin{itemize}
    \item \textbf{Blocked Views from Thick Blood or Cataracts:} Even though near-infrared light passes through thin blood better than visible light, thick bleeding under the retina, dense cataracts, or blood floating in the eye gel can still hide underlying vessel problems.
    \item \textbf{Overlapping Glow Patterns:} Harmless fluid pockets, tiny tears in the deep retinal layer, or chronic leaky vessels can glow late in the test and look very similar to hidden abnormal blood vessels.
    \item \textbf{Faint Signal or Missed Timing:} Because the infrared glow is faint, blinking, trouble holding focus, or missing the very first seconds of dye filling can make it harder to spot fast-filling vessels or tumor patterns.
    \item \textbf{Combining Tests for Clear Answers:} Eye specialists almost always read ICGA alongside cross-sectional OCT scans, OCT angiography, and fluorescein angiography to confirm what is happening. If thick blood blocks the view, the doctor may repeat targeted imaging once the blood begins to clear.
\end{itemize}

\section{Augmenting Methods}

\begin{itemize}
    \item \textbf{Simultaneous Fluorescein Angiography (FFA):} Dual-mode cameras can capture both fluorescein and ICG images at the same time, combining sharp views of surface retinal vessels with deep choroidal maps.
    \item \textbf{Optical Coherence Tomography (OCT):} Painless cross-sectional OCT scans measure the thickness of the deep choroid layer, show fluid pockets beneath the retina, and track how well treatments are working over time.
    \item \textbf{OCT Angiography (OCTA):} A dye-free scan that maps blood flow beneath the retina, making it easy to monitor abnormal vessel networks at follow-up visits without repeating IV injections.
    \item \textbf{Autofluorescence and Eye Ultrasound:} Specialized light photos (fundus autofluorescence) reveal old fluid tracks, while gentle eye ultrasound and blood tests help evaluate eye tumors or body-wide inflammatory conditions.
\end{itemize}

\section{Comparison with Alternative Modalities}

\begin{itemize}
    \item \textbf{Compared with Fluorescein Angiography (FFA):} Fluorescein angiography is better for spotting leaks in fine surface retinal vessels, while near-infrared ICGA sees much deeper through pigment and thin blood, making it the gold standard for polypoidal choroidal vasculopathy and deep choroidal inflammation.
    \item \textbf{Compared with OCT Angiography (OCTA):} OCTA needs no dye injection and maps branching vessel networks quickly, while ICGA is better at showing slow-flowing vessel bulges (polyps), widespread vessel leakiness, and deep inflammatory spots.
    \item \textbf{Compared with Enhanced-Depth OCT:} Cross-sectional OCT is fast, dye-free, and ideal for tracking fluid and choroid thickness at every visit, while ICGA provides a wide overhead map of blood flow and leakiness to guide targeted laser treatments.
\end{itemize}

\section{Follow-up Steps}

\begin{itemize}
    \item \textbf{Short Post-Test Check and Thyroid-Test Reminder:} Staff will briefly watch you for any delayed allergy signs and check your IV site. Remember to tell your other doctors that you had ICG dye if you have a radioactive iodine thyroid test scheduled within the next week.
    \item \textbf{Eye-Drop Safety After You Leave:} Wear sunglasses outdoors and wait until your blurry near vision and glare sensitivity fade (usually 4 to 6 hours) before driving. Seek urgent eye care if you develop severe eye pain, headache, rainbow halos around lights, or nausea, which can signal a sudden spike in eye pressure.
    \item \textbf{Personalized Treatment Plan:} Your retina specialist combines your ICGA maps with OCT and other eye scans to plan targeted care---such as eye injections (anti-VEGF), gentle photodynamic laser therapy, or anti-inflammatory medicines.
    \item \textbf{Ongoing Monitoring:} Most follow-up visits use quick, dye-free OCT scans to track fluid and healing. Repeat ICGA is usually reserved for new symptoms, suspected return of abnormal vessel bulges, or poor response to treatment.
\end{itemize}


\section{What You Should Know About Indocyanine Green Angiography?}

\begin{itemize}
    \item \textbf{Gentle Lighting and Minimal Discomfort:} Aside from a quick IV pinch in your arm, the test is painless and nothing touches your eye. The dim red near-infrared light is much gentler than the bright flashes of standard fluorescein eye photography, and nausea is uncommon.
    \item \textbf{Zero X-Ray Radiation and Dye Safety:} ICGA uses safe near-infrared light instead of X-rays. Because the dye contains a small amount of sodium iodide and is cleared by the liver, rare allergic reactions can happen, so clinics always keep emergency treatment close by.
    \item \textbf{Visit Length and Results:} Plan for about 45 to 60 minutes in the clinic, including eye dilation and relaxed breaks between the timed photos. Your eye specialist often reviews the findings with you the same day or at your next visit.
    \item \textbf{Why Your Doctor Orders It:} Near-infrared light sees deeper behind the retina and through thin layers of blood than standard tests, helping diagnose hidden deep vessel problems like polypoidal choroidal vasculopathy or chronic fluid buildup.
    \item \textbf{What to Tell Your Care Team:} Always share any history of iodide or dye allergy, overactive thyroid, liver disease, pregnancy, or upcoming thyroid scans beforehand. Bring sunglasses and a driver home, and tell staff right away if you feel itchy, flushed, or short of breath.
\end{itemize}

\section{Key References}

\begin{itemize}
    \item Yannuzzi LA. Indocyanine Green Angiography: A Perspective on Use in the Clinical Setting. \textit{American Journal of Ophthalmology}, 2011;151(5):745--751.
    \item American Academy of Ophthalmology (AAO). \textit{Basic and Clinical Science Course (BCSC), Section 12: Retina and Vitreous}. San Francisco, CA: AAO, 2023--2024.
    \item Cheung CMG, Lai TYY, Ruamviboonsuk P, et al. Polypoidal Choroidal Vasculopathy: Definition, Pathogenesis, Diagnosis, and Management. \textit{Ophthalmology}, 2018;125(5):708--724.
    \item Hope-Ross M, Yannuzzi LA, Gragoudas ES, et al. Adverse Reactions Due to Indocyanine Green. \textit{Ophthalmology}, 1994;101(3):529--533.
\end{itemize}
